% Bagian: first-order-logic
% Bab: syntax-and-semantics
% Subbagian: assignments

\documentclass[../../../include/open-logic-section]{subfiles}

\begin{document}

\olfileid[id]{fol}{syn}{ass}

\olsection{Penugasan Variabel}

\begin{explain}
Suatu penugasan !!{variable}~$s$ memberikan nilai kepada \emph{setiap}
variabel---dan jumlahnya tak hingga. Tentu saja, hal ini tidak diperlukan.
Kita mengharuskan penugasan !!{variable} menetapkan nilai kepada semua
!!{variable}s semata-mata karena hal itu sangat mempermudah pembahasan.
Nilai suatu suku~$t$, dan apakah suatu !!a{formula}~$!A$ terpenuhi atau
tidak dalam !!a{structure} sehubungan dengan~$s$, hanya bergantung pada
penetapan yang dibuat~$s$ kepada !!{variable}s dalam~$t$ dan kepada
!!{variable}s bebas dari~$!A$. Inilah isi dua proposisi berikutnya. Untuk
membuat gagasan ``bergantung pada'' tepat, kita menunjukkan bahwa setiap dua
penugasan variabel yang bersepakat pada semua variabel dalam~$t$ memberikan
nilai yang sama, dan bahwa $!A$ terpenuhi relatif terhadap penugasan yang
satu jika dan hanya jika terpenuhi relatif terhadap yang lain apabila kedua
penugasan variabel tersebut bersepakat pada semua variabel bebas dari~$!A$.
\end{explain}

\begin{prop}\ollabel{prop:valindep}
Jika !!{variable}s dalam suatu suku~$t$ terdapat di antara $x_1$, \dots,
~$x_n$, dan $s_1(x_i) = s_2(x_i)$ untuk $i = 1$, \dots,~$n$, maka
$\Value{t}{M}[s_1] = \Value{t}{M}[s_2]$.
\end{prop}

\begin{proof}
Dengan induksi pada kompleksitas~$t$. Bagi kasus dasar, $t$ dapat berupa
!!a{constant} atau salah satu variabel~$x_1$, \dots,~$x_n$. Jika $t = c$,
maka $\Value{t}{M}[s_1] = \Assign{c}{M} = \Value{t}{M}[s_2]$. Jika
$t = x_i$, maka $s_1(x_i) = s_2(x_i)$ berdasarkan hipotesis proposisi,
sehingga $\Value{t}{M}[s_1] = s_1(x_i) = s_2(x_i) =
\Value{t}{M}[s_2]$.

Bagi langkah induktif, andaikan $t = \Atom{f}{t_1, \dots, t_k}$ dan bahwa
klaim berlaku bagi $t_1$, \dots, $t_k$. Maka
\begin{align*}
  \Value{t}{M}[s_1] & = \Value{\Atom{f}{t_1, \dots, t_k}}{M}[s_1] \\
  & = \Assign{f}{M}(\Value{t_1}{M}[s_1], \dots, \Value{t_k}{M}[s_1]).
\intertext{Untuk $j = 1$, \dots,~$k$, !!{variable}s dari~$t_j$ terdapat
    di antara $x_1$, \dots,~$x_n$. Berdasarkan hipotesis induksi,
    $\Value{t_j}{M}[s_1] = \Value{t_j}{M}[s_2]$. Jadi,}
\Value{t}{M}[s_1] & = \Value{\Atom{f}{t_1, \dots, t_k}}{M}[s_1] \\
  & = \Assign{f}{M}(\Value{t_1}{M}[s_1], \dots, \Value{t_k}{M}[s_1]) \\
  & = \Assign{f}{M}(\Value{t_1}{M}[s_2], \dots, \Value{t_k}{M}[s_2]) \\
  & = \Value{\Atom{f}{t_1, \dots, t_k}}{M}[s_2] = \Value{t}{M}[s_2].
\end{align*}
\end{proof}

\begin{prop}\ollabel{prop:satindep}
Jika !!{variable}s bebas dalam $!A$ terdapat di antara $x_1$, \dots,~$x_n$,
dan $s_1(x_i) = s_2(x_i)$ untuk $i = 1$, \dots,~$n$, maka
$\Sat{M}{!A}[s_1]$ jika dan hanya jika $\Sat{M}{!A}[s_2]$.
\end{prop}

\begin{proof}
Kita menggunakan induksi pada kompleksitas $!A$. Bagi kasus dasar, ketika
$!A$ atomik, $!A$ dapat berupa:
\iftag{prvTrue}{$\ltrue$,}{}
\iftag{prvFalse}{$\lfalse$,}{}
$\Atom{R}{t_1, \dots, t_k}$ bagi suatu predikat $k$-tempat $R$ dan suku
$t_1$, \dots,~$t_k$, atau $\eq[t_1][t_2]$ bagi suku $t_1$ dan~$t_2$.
Dalam dua kasus terakhir, kita hanya menunjukkan arah maju dari
!!{biconditional}, sebab pembuktian arah sebaliknya simetris.

\begin{enumerate}
\tagitem{prvTrue}{%
  \indcase{!A}{\ltrue}{baik $\Sat{M}{\indfrm}[s_1]$ maupun
    $\Sat{M}{\indfrm}[s_2]$.}}{}

\tagitem{prvFalse}{%
  \indcase{!A}{\lfalse}{baik $\Sat/{M}{!A}[s_1]$ maupun
    $\Sat/{M}{!A}[s_2]$.}}{}

\item
  \indcase{!A}{\Atom{R}{t_1, \ldots, t_k}}{andaikan
    $\Sat{M}{\indfrm}[s_1]$. Maka
    \[
    \langle \Value{t_1}{M}[s_1], \ldots, \Value{t_k}{M}[s_1] \rangle
    \in \Assign{R}{M}.
    \]
    Untuk $i = 1$, \dots,~$k$, $\Value{t_i}{M}[s_1] =
    \Value{t_i}{M}[s_2]$ berdasarkan \olref{prop:valindep}. Jadi, kita juga
    mempunyai $\langle \Value{t_1}{M}[s_2], \ldots,
    \Value{t_k}{M}[s_2] \rangle \in \Assign{R}{M}$, dan karenanya
    $\Sat{M}{\indfrm}[s_2]$.}
\item
  \indcase{!A}{\eq[t_1][t_2]}{andaikan $\Sat{M}{\indfrm}[s_1]$.
    Maka $\Value{t_1}{M}[s_1] = \Value{t_2}{M}[s_1]$. Jadi,
    \begin{align*}
      \Value{t_1}{M}[s_2] & = \Value{t_1}{M}[s_1]
      & \text{(berdasarkan \olref{prop:valindep})} \\
      & = \Value{t_2}{M}[s_1]
      & \text{(karena $\Sat{M}{\eq[t_1][t_2]}[s_1]$)}\\
      &= \Value{t_2}{M}[s_2]
      & \text{(berdasarkan \olref{prop:valindep}),}
    \end{align*}
    sehingga $\Sat{M}{\eq[t_1][t_2]}[s_2]$.}
\end{enumerate}

Sekarang andaikan $\Sat{M}{!B}[s_1]$ jika dan hanya jika
$\Sat{M}{!B}[s_2]$ bagi semua !!{formula}s $!B$ yang kurang kompleks
daripada~$!A$. Langkah induksi berjalan menurut kasus-kasus yang ditentukan
oleh operator utama dari~$!A$. Dalam setiap kasus, kita hanya menunjukkan
arah maju dari !!{biconditional}; pembuktian arah sebaliknya simetris. Dalam
semua kasus selain kasus kuantor, kita menerapkan hipotesis induksi kepada
!!{subformula}s~$!B$ dari~$!A$. Variabel bebas dari~$!B$ terdapat di antara
variabel bebas dari~$!A$. Dengan demikian, jika $s_1$ dan $s_2$ bersepakat
pada variabel bebas dari~$!A$, keduanya juga bersepakat pada variabel bebas
dari~$!B$, dan hipotesis induksi berlaku bagi~$!B$.

\begin{enumerate}
\tagitem{defNot}{}{%
  \iftag{probNot}{%
    \indcase!{!A}{\lnot !B}{}}{%
    \indcase{!A}{\lnot !B}{jika $\Sat{M}{\indfrm}[s_1]$, maka
      $\Sat/{M}{!B}[s_1]$, sehingga berdasarkan hipotesis induksi,
      $\Sat/{M}{!B}[s_2]$, dan karenanya $\Sat{M}{\indfrm}[s_2]$.}}}

\tagitem{defAnd}{}{%
  \iftag{probAnd}{%
    \indcase!{!A}{!B \land !C}{}}{%
    \indcase{!A}{!B \land !C}{jika $\Sat{M}{\indfrm}[s_1]$, maka
      $\Sat{M}{!B}[s_1]$ dan $\Sat{M}{!C}[s_1]$, sehingga berdasarkan
      hipotesis induksi, $\Sat{M}{!B}[s_2]$ dan $\Sat{M}{!C}[s_2]$.
      Dengan demikian, $\Sat{M}{\indfrm}[s_2]$.}}}

\tagitem{defOr}{}{%
  \iftag{probOr}{%
    \indcase!{!A}{!B \lor !C}{}}{%
    \indcase{!A}{!B \lor !C}{jika $\Sat{M}{\indfrm}[s_1]$, maka
      $\Sat{M}{!B}[s_1]$ atau $\Sat{M}{!C}[s_1]$. Berdasarkan hipotesis
      induksi, $\Sat{M}{!B}[s_2]$ atau $\Sat{M}{!C}[s_2]$, sehingga
      $\Sat{M}{\indfrm}[s_2]$.}}}

\tagitem{defIf}{}{%
  \iftag{probIf}{%
    \indcase!{!A}{!B \lif !C}{}}{%
    \indcase{!A}{!B \lif !C}{jika $\Sat{M}{\indfrm}[s_1]$, maka
    $\Sat/{M}{!B}[s_1]$ atau $\Sat{M}{!C}[s_1]$. Berdasarkan hipotesis
    induksi, $\Sat/{M}{!B}[s_2]$ atau $\Sat{M}{!C}[s_2]$, sehingga
    $\Sat{M}{\indfrm}[s_2]$.}}}

\tagitem{defIff}{}{%
  \iftag{probIff}{%
    \indcase!{!A}{!B \liff !C}{}}{%
    \indcase{!A}{!B \liff !C}{jika $\Sat{M}{\indfrm}[s_1]$, maka
      $\Sat{M}{!B}[s_1]$ dan $\Sat{M}{!C}[s_1]$, atau
      $\Sat/{M}{!B}[s_1]$ dan $\Sat/{M}{!C}[s_1]$. Berdasarkan hipotesis
      induksi, $\Sat{M}{!B}[s_2]$ dan $\Sat{M}{!C}[s_2]$, atau
      $\Sat/{M}{!B}[s_2]$ dan $\Sat/{M}{!C}[s_2]$. Dalam kedua kasus,
      $\Sat{M}{\indfrm}[s_2]$.}}}

\tagitem{defEx}{}{%
  \iftag{probEx}{%
    \indcase!{!A}{\lexists[x][!B]}{}}{%
    \indcase{!A}{\lexists[x][!B]}{jika $\Sat{M}{\indfrm}[s_1]$, terdapat
      $m \in \Domain{M}$ sedemikian sehingga
      $\Sat{M}{!B}[\Subst{s_1}{m}{x}]$. Misalkan
      $s_1' = \Subst{s_1}{m}{x}$ dan $s_2' =
      \Subst{s_2}{m}{x}$. Variabel bebas dari~$!B$ terdapat di antara
      $x_1$, \dots, $x_n$, dan $x$. $s_1'(x_i) = s_2'(x_i)$, sebab $s_1'$
      dan $s_2'$ masing-masing merupakan varian-$x$ dari $s_1$ dan~$s_2$,
      dan berdasarkan hipotesis $s_1(x_i) = s_2(x_i)$.
      $s_1'(x) = s_2'(x) = m$ berdasarkan cara kita mendefinisikan
      $s_1'$ dan~$s_2'$. Maka
      hipotesis induksi berlaku bagi $!B$ dan $s_1'$, $s_2'$, sehingga
      $\Sat{M}{!B}[s_2']$. Dengan demikian, karena
      $s_2' = \Subst{s_2}{m}{x}$, terdapat $m \in \Domain{M}$ sedemikian
      sehingga $\Sat{M}{!B}[\Subst{s_2}{m}{x}]$, dan karenanya
      $\Sat{M}{\indfrm}[s_2]$.}}}

\tagitem{defAll}{}{%
  \iftag{probAll}{%
    \indcase!{!A}{\lforall[x][!B]}{}}{%
    \indcase{!A}{\lforall[x][!B]}{jika $\Sat{M}{\indfrm}[s_1]$, maka bagi
      setiap $m \in \Domain{M}$, $\Sat{M}{!B}[\Subst{s_1}{m}{x}]$. Kita
      ingin menunjukkan bahwa bagi setiap $m \in \Domain{M}$ juga berlaku
      $\Sat{M}{!B}[\Subst{s_2}{m}{x}]$. Jadi, misalkan $m \in \Domain{M}$
      sebarang, dan perhatikan $s_1' = \Subst{s_1}{m}{x}$ serta $s_2' =
      \Subst{s_2}{m}{x}$. Kita mempunyai $\Sat{M}{!B}[s_1']$. Variabel bebas
      dari~$!B$ terdapat di antara $x_1$, \dots, $x_n$, dan $x$.
      $s_1'(x_i) = s_2'(x_i)$, sebab $s_1'$ dan $s_2'$ masing-masing
      merupakan varian-$x$ dari $s_1$ dan~$s_2$, dan berdasarkan hipotesis
      $s_1(x_i) = s_2(x_i)$. Kita mempunyai
      $s_1'(x) = s_2'(x) = m$ berdasarkan cara kita mendefinisikan
      $s_1'$ dan~$s_2'$. Maka hipotesis induksi berlaku
      bagi~$!B$ serta~$s_1'$, $s_2'$, dan kita mempunyai
      $\Sat{M}{!B}[s_2']$. Hal ini berlaku bagi setiap $m \in \Domain{M}$,
      yakni, $\Sat{M}{!B}[\Subst{s_2}{m}{x}]$ bagi semua~$m \in
      \Domain{M}$, sehingga $\Sat{M}{\indfrm}[s_2]$.}}}
\end{enumerate}
Dengan induksi, kita memperoleh bahwa $\Sat{M}{!A}[s_1]$ jika dan hanya jika
$\Sat{M}{!A}[s_2]$ setiap kali !!{variable}s bebas dalam $!A$ terdapat di
antara $x_1$, \dots, $x_n$ dan $s_1(x_i)=s_2(x_i)$ untuk
$i = 1$, \dots,~$n$.
\end{proof}

\begin{probtag}{probNot,probOr,probAnd,probIf,probIff,probEx,probAll}
Lengkapilah pembuktian \olref[fol][syn][ass]{prop:satindep}.
\end{probtag}

\begin{explain}
!!^{sentence}s tidak mempunyai variabel bebas, sehingga setiap dua
penugasan variabel menetapkan hal yang sama kepada semua variabel bebas
(yang jumlahnya nol) dari setiap kalimat. Dengan demikian, proposisi yang
baru saja dibuktikan berarti bahwa apakah !!a{sentence} terpenuhi atau tidak
dalam suatu struktur relatif terhadap penugasan variabel sama sekali tidak
bergantung pada penugasan tersebut. Kita akan mencatat fakta ini. Fakta ini
membenarkan definisi keterpenuhan !!a{sentence} dalam !!a{structure} (tanpa
menyebutkan penugasan variabel) berikut ini.
\end{explain}

\begin{cor}
\ollabel{cor:sat-sentence}
Jika $!A$ merupakan !!a{sentence} dan $s$ suatu penugasan variabel, maka
$\Sat{M}{!A}[s]$ jika dan hanya jika $\Sat{M}{!A}[s']$ bagi setiap
penugasan variabel~$s'$.
\end{cor}

\begin{proof}
  Misalkan $s'$ sebarang penugasan variabel. Karena $!A$ merupakan
  !!a{sentence}, kalimat itu tidak mempunyai variabel bebas, sehingga setiap
  penugasan variabel~$s'$ secara trivial menetapkan hal yang sama kepada
  semua variabel bebas dari~$!A$ seperti yang ditetapkan~$s$. Jadi, syarat
  \olref{prop:satindep} terpenuhi, dan kita mempunyai
  $\Sat{M}{!A}[s]$ jika dan hanya jika $\Sat{M}{!A}[s']$.
\end{proof}

\begin{defn}
\ollabel{defn:satisfaction}
Jika $!A$ merupakan !!a{sentence}, kita mengatakan bahwa
!!a{structure}~$\Struct M$ \emph{memenuhi}~$!A$, $\Sat{M}{!A}$, jika dan
hanya jika $\Sat{M}{!A}[s]$ bagi semua penugasan variabel~$s$.
\end{defn}

Jika $\Sat{M}{!A}$, kita juga cukup mengatakan bahwa \emph{$!A$ benar
dalam~$\Struct{M}$.} Gagasan keterpenuhan secara alamiah diperluas dari
!!{sentence}s individual ke himpunan !!{sentence}s.

\begin{defn}
\ollabel{defn:sat}
Jika $\Gamma$ merupakan himpunan !!{sentence}s, kita mengatakan
bahwa !!a{structure}~$\Struct M$ \emph{memenuhi}~$\Gamma$,
$\Sat{M}{\Gamma}$, jika dan hanya jika $\Sat{M}{!A}$ bagi semua
$!A \in \Gamma$.
\end{defn}

\begin{prop}\ollabel{prop:sentence-sat-true}
  Misalkan $\Struct{M}$ suatu !!a{structure}, $!A$ suatu !!a{sentence}, dan
  $s$ suatu penugasan variabel. $\Sat{M}{!A}$ jika dan hanya jika
  $\Sat{M}{!A}[s]$.
\end{prop}

\begin{proof}
Latihan.
\end{proof}

\begin{prob}
Buktikan \olref[fol][syn][ass]{prop:sentence-sat-true}
\end{prob}

\begin{prop}\ollabel{prop:sat-quant}
  Andaikan $!A(x)$ hanya memuat $x$ secara bebas, dan $\Struct M$ merupakan
  !!a{structure}. Maka:
  \begin{tagenumerate}{prvEx,prvAll}
    \tagitem{prvEx}{$\Sat{M}{\lexists[x][!A(x)]}$ jika dan hanya jika
      $\Sat{M}{!A(x)}[s]$ bagi setidaknya satu penugasan variabel~$s$.}{}
    \tagitem{prvAll}{$\Sat{M}{\lforall[x][!A(x)]}$ jika dan hanya jika
      $\Sat{M}{!A(x)}[s]$ bagi semua penugasan variabel~$s$.}{}
\end{tagenumerate}
\end{prop}

\begin{proof}
Latihan.
\end{proof}

\begin{prob}
Buktikan \olref[fol][syn][ass]{prop:sat-quant}.
\end{prob}

\begin{prob}
\DeclareRobustCommand{\VDash}{\mathrel{||}\joinrel\Relbar}
Andaikan $\Lang L$ merupakan bahasa tanpa !!{function}s. Misalkan diberikan
!!{structure}~$\Struct M$, suatu !!a{constant}~$c$, dan
$a \in \Domain M$. Definisikan $\Struct M[a/c]$ sebagai
!!{structure} yang sama persis dengan~$\Struct M$, kecuali bahwa
$\Assign{c}{M[a/c]} = a$. Definisikan $\Struct M \VDash !A$ bagi
!!{sentence}s~$!A$ dengan:
\begin{enumerate}
\tagitem{prvFalse}{%
  \indcase{!A}{\lfalse}{tidak berlaku $\Struct{M} \VDash \indfrm$.}}{}

\tagitem{prvTrue}{%
  \indcase{!A}{\ltrue}{$\Struct{M} \VDash \indfrm$.}}{}

\item \indcase{!A}{\Atom{R}{d_1, \dots, d_n}}{$\Struct M \VDash \indfrm$
  jika dan hanya jika
  $\langle \Assign{d_1}{M}, \dots, \Assign{d_n}{M} \rangle \in
  \Assign{R}{M}$.}

\item \indcase{!A}{\eq[d_1][d_2]}{$\Struct M \VDash \indfrm$ jika dan hanya
  jika $\Assign{d_1}{M} = \Assign{d_2}{M}$.}

\tagitem{prvNot}{%
  \indcase{!A}{\lnot !B}{$\Struct{M} \VDash \indfrm$ jika dan hanya jika
    tidak berlaku $\Struct{M} \VDash {!B}$.}}{}

\tagitem{prvAnd}{%
  \indcase{!A}{(!B \land !C)}{$\Struct{M} \VDash
    \indfrm$ jika dan hanya jika $\Struct{M} \VDash!B$ dan
    $\Struct{M} \VDash !C$.}}{}

\tagitem{prvOr}{%
  \indcase{!A}{(!B \lor !C)}{$\Struct{M} \VDash \indfrm$ jika dan hanya jika
    $\Struct{M} \VDash !B$ atau $\Struct{M} \VDash !C$ (atau keduanya).}}{}

\tagitem{prvIf}{%
  \indcase{!A}{(!B \lif !C)}{$\Struct{M} \VDash \indfrm$ jika dan hanya jika
    tidak berlaku $\Struct {M} \VDash !B$ atau $\Struct M \VDash !C$
    (atau keduanya).}}{}

\tagitem{prvIff}{%
  \indcase{!A}{(!B \liff !C)}{$\Struct{M} \VDash {\indfrm}$ jika dan hanya
    jika baik $\Struct{M} \VDash {!B}$ maupun $\Struct{M} \VDash {!C}$
    keduanya berlaku, atau baik $\Struct{M} \VDash {!B}$ maupun
    $\Struct{M} \VDash {!C}$ keduanya tidak berlaku.}}{}

\tagitem{prvAll}{%
  \indcase{!A}{\lforall[x][!B]}{$\Struct{M} \VDash {\indfrm}$ jika dan hanya
    jika bagi semua $a \in \Domain{M}$,
    $\Struct{M[a/c]} \VDash \Subst{!B}{c}{x}$, jika $c$ tidak muncul
    dalam~$!B$.}}{}

\tagitem{prvEx}{%
  \indcase{!A}{\lexists[x][!B]}{$\Struct{M} \VDash
  {\indfrm}$ jika dan hanya jika terdapat $a \in \Domain M$ sedemikian
  sehingga $\Struct{M[a/c]} \VDash \Subst{!B}{c}{x}$, jika $c$ tidak
  muncul dalam~$!B$.}}{}
\end{enumerate}
Misalkan $x_1$, \dots, $x_n$ semua !!{variable}s bebas dalam~$!A$,
$c_1$, \dots, $c_n$ simbol konstanta yang tidak terdapat dalam~$!A$,
$a_1$, \dots, $a_n \in \Domain M$, dan $s(x_i) = a_i$.

Tunjukkan bahwa $\Sat{M}{!A}[s]$ jika dan hanya jika
$\Struct M[a_1/c_1,\dots,a_n/c_n]
\VDash \Subst{\Subst{!A}{c_1}{x_1}\dots}{c_n}{x_n}$.

(Soal ini menunjukkan bahwa semantik bagi logika orde pertama dapat
diberikan tanpa menggunakan penugasan variabel.)
\end{prob}

\begin{prob}
Andaikan $f$ merupakan simbol fungsi yang tidak terdapat dalam~$!A(x,y)$.
Tunjukkan bahwa terdapat !!a{structure}~$\Struct{M}$ sedemikian sehingga
$\Sat{M}{\lforall[x][\lexists[y][!A(x,y)]]}$ jika dan hanya jika terdapat
suatu~$\Struct M'$ sedemikian sehingga
$\Sat{M'}{\lforall[x][!A(x,f(x))]}$.

(Soal ini merupakan kasus khusus dari apa yang dikenal sebagai Teorema
Skolem; $\lforall[x][!A(x,f(x))]$ disebut suatu \emph{bentuk normal Skolem}
dari $\lforall[x][\lexists[y][!A(x,y)]]$.)
\end{prob}

\end{document}
